\section*{Limitations}
\label{sec:limitations}

\textbf{Guarantees.} Proposition~\ref{prop:termination} is a bounded termination-and-escalation property, not a convergence theorem. Calibration controls evidence omission under its calibration distribution, not correct legal conclusions, temporal-shift robustness, or freedom from LLM hallucinations~\citep{dahl2024hallucinating}. Because litigation streams are non-stationary, the omission guarantee is weakest under the temporal split (Table~\ref{tab:risk_control_temporal})---precisely where deployment needs it most---so we treat rolling or weighted re-calibration as a deployment requirement rather than an optional refinement.

\textbf{Unfinished external-generality check.} The central mechanism is a domain-general validity-gating operator, but its transfer to \emph{external} retrieval stacks is specified only as a pre-registered protocol (Protocol~P1, Appendix~\ref{app:p1}) and is not yet executed; the accuracy ordering is shown to persist only across our own backbones with an open-weights stack (Table~\ref{tab:robustness}). Cross-stack transfer therefore remains the main open generality test.

\textbf{Benchmark scale and evaluation.} At the current scale PLRE remains an initial benchmark over five product archetypes and 312 product-mapping cases, so it should be read as evidence for the validity-gated entailment formulation rather than as coverage of all AI-product litigation. The report-quality study (RQ2) uses two blinded raters and is treated as supporting, not primary, evidence. Because 847 of 1{,}247 labels are retained pre-labels, all primary significance claims are restricted to the human-adjudicated subset, and effect sizes over the strongest baseline are modest (Table~\ref{tab:human_only_main}).

\textbf{Coverage and jurisdiction.} The dataset is skewed toward the largest defendants---the top three account for 41.3\% of cases---limiting representativeness for smaller vendors. ALERT targets U.S.\ federal CourtListener/RECAP coverage and English-language documents; non-U.S.\ jurisdictions require different docket schemas and citation practices, and extending the treatment graph to them is future work.

\textbf{Extractor recall.} Negative-treatment edge recall (0.77 internal gold; 20.4\% missed-closure rate against an authoritative citator) is useful but not lossless, and a missed closure is exactly what allows a stale precedent to leak into a risk decision; ALERT mitigates this with human-review routing for low-confidence closures, but improving extraction recall remains an open problem.
